\section{问题重述}

\subsection{问题背景}

微网是为解决光伏发电等分布式能源本地消纳而快速发展起来的小型发用电系统，由分布式光伏、
储能设备与小区负载构成，并可通过联络线与外部电网（下称外网）交换电能。由于外网电价在
不同时段存在显著差异，非孤岛式微网可在电价较低或分布式能源有剩余电量的时段为储能充电，
在电价较高的时段释放电能，从而在满足小区负载的前提下降低购电费用。

微网与外网调控的核心任务是：利用小区负载、光伏发电量与储能设备当前储电量等数据，
尽可能精准地给出微网的购电量与储能设备的充放电量，使微网提供的电力既满足小区负载，
又尽可能节省购电费用。

\subsection{需要解决的问题}

\begin{enumerate}[label=\textbf{问题\arabic*.}, leftmargin=4em, itemsep=2pt]
\item \textbf{单典型日最优计划}：每天电价与小区负载相同、光伏发电功率有一天的预测数据，
要求微网提供的电能不低于小区负载，且储能设备在 0:00 与 24:00 的储电量相同。
建立每天 0:00 制定当天计划购电策略的数学模型，给出指定时段与全天的购电量、购电费，
以及储能各时段的充放电量与 0:00、24:00 储电量。
\item \textbf{全年逐日计划与紧急购电}：电价每天相同，但小区负载与光伏发电功率逐日变化；
若微网提供的电能低于负载，需按交易时刻电价的 \textbf{5 倍}紧急购电，其余时段购电费按计划购电量计。
在每天 0:00 制定当天的计划购电策略，覆盖 2025 年 2 月 1 日至 12 月 31 日。
\item \textbf{滚动预报下的计划与调整}：每天 0:00、6:00、12:00、18:00 可获得未来 24 小时整点的
光伏发电功率预报，可据 0:00 预报制定计划、并据其他时刻预报调整策略。计划购电量高于调整购电量的部分，
违约电价为交易时刻电价的 \textbf{50\%}；调整购电量高于计划购电量的部分，超出部分电价为
交易时刻电价的 \textbf{1.5 倍}；总购电费用包括计划购电费用、紧急购电费用与调整购电量相关费用。
要求建立模型并分析是否需要引入其他时刻的预报制定调整策略。
\item \textbf{实时波动电价}：外网电价实际是实时波动的，需针对波动电价建立当天购电策略模型，
并在波动电价下重新计算问题二与问题三。
\end{enumerate}

\subsection{数据与参数说明}

\begin{table}[htbp]
\centering
\caption{原始数据与储能参数清单}
\label{tab:data}
\small
\begin{tabularx}{\textwidth}{@{}p{1.5cm}p{2.4cm}p{6.2cm}X@{}}
\toprule
来源 & 内容 & 结构 & 单位与口径 \\
\midrule
附件1 & 电价、小区负载、光伏预测功率 & 单典型日 144 个 10 分钟时段（列头为区间结束时刻，如表头 \texttt{00:10} 表示 0:00--0:10） & 元/kWh；kW；kW \\
附录1 & 储能参数 & 最大容量 12000 kWh；最大充放电功率 5000 kW；SOC 限幅 1200--10800 kWh；2025-01-01 0:00 储电量 6000 kWh；充放电效率 90\% & kWh、kW \\
附件2 & 小区负载、光伏发电实际功率 & 2025-01-01 至 12-31，365 天 $\times$ 144 时段（两张工作表） & kW \\
附件3 & 光伏发电预报 & 365 天 $\times$ 4 个发布时刻（0:00/6:00/12:00/18:00）$\times$ 未来 24 个整点 & kW \\
附件4 & 实时电价 & 2025-01-01 至 12-31，365 天 $\times$ 144 时段 & 元/kWh \\
附件5 & 结果模板 & result1/2/3/4-2/4-3.xlsx & kWh、元 \\
\bottomrule
\end{tabularx}
\end{table}

数据核验表明：附件1/2/4 无缺失值、无负值；附件3 的 1095 个空值全部为每日第 2--4 行
“日期”列的合并单元格结构空值（向下填充即可），24 个数值列无缺失。
